\head{\textsc{\emph{Ilíada}, Canto I}}
\section*{\centering\textbf{\textsc{Canto} I}\footnote{\neh, pp. 58-80.}}
\addcontentsline{toc}{subsection}{\textsc{Canto I}}
%
\beginnumbering
\pstart
\lettrine[lines=4, nindent=2pt]{C}{\textsc{anta, diosa, el encono}} de Aquiles Peleyades,\\
que, aciago, a los aqueos causó males sin cuento,\\
y tantas nobles almas de héroes echó al Hades,\\
dando a perros y buitres sus cuerpos de alimento\\
(pues así los designios Zeus se realizaron)\\
desde que por primera vez riñendo, quebraron\\
el Atrida, rey de hombres, y el deífico Aquiles.\par
¿Cuál fué de entre los dioses el que los puso hostiles?\\
El hijo de Latona y Zeus, quien resentido\\
con el rey, al ejército envió maligna peste\\
que a las tropas diezmaba, porque había ofendido\\
Atrida al sacerdote Crises, cuando llegó éste\\
a las veleras naves aqueas, aportando\\
por rescate de su hija multa inmensa, y llevando\\
en su mano las ínfulas del flechador Apolo\\
a un cetro de oro atadas; con que imploró, no sólo\\
a ambos Atridas, jefes de pueblos, sino a cada\\
\dlgo{uno de los aqueos, diciéndoles:}{--Oh Atridas}\\
y soldados aqueos de las grebas pulidas,\\
quieran los dioses desde su olímpica morada\\
que la ciudad de Príamo toméis, y con buen sino\\
volváis a vuestra tierra; mas, dadme a mi hija amada,\\
aceptando el rescate, y acatando al divino\\
\dlgo{flechero Apolo.}{Todos aclaman que se deba}\\
respeto al sacerdote, recibiendo la multa.\\
Agamenón Atrida solamente no aprueba;\\
antes, con malos modos, lo echa de allá y lo insulta:\par
--Que no vuelva yo, anciano, ni ahora ni luego a hallarte,\\
ya te quedes o tornes, junto a la flota anclada;\\
pues ni el cetro ni la ínfula del dios podrán guardarte.\\
No soltaré a tu hija, antes que envejezca expatriada,\\
tejiendo en casa, en Argos, y durmiendo conmigo.\\
Y vete, y no me irrites, y evita así el castigo.\par
Dice, y medroso el viejo, lo acata y abandona\\
en silencio, y costeando va el estruendoso mar;\\
y al rey Apolo, vástago de la crespa Latona,\\
mientras se aparta, empieza con fervor a rogar:\par
--Oyeme, dios del arco de plata, protector\\
de Crisa y de la santa Kila, dominador\\
de Ténedos, ¡oh Esmintio!: Si tu iglesia graciosa\\
labrar supe, y quemarte pingües muslos de cabras\\
y toros, que los dánaos, conforme a estas palabras,\\
mis lágrimas expíen con tu flecha luctuosa.\par
Rogó así, y Febo Apolo que lo oyó, con airado\\
corazón, baja al punto del Olimpo encumbrado,\\
al hombro arco y aljaba de bien cerrado broche.\\
Y en la espalda resuénanle, al moverse enojado,\\
las flechas, cuando avanza semejante a la noche.\\
Siéntese después, lejos de las naves, y tira\\
una flecha, al terrible son del arco de plata.\\
Mulas y ágiles perros hiere primero su ira;\\
mas, luego, con fatales tiros los hombres mata,\\
y arden sin tregua innúmeros difuntos en la pira.\par
Nueve días las tropas flechó el dios por doquiera;\\
mas, el décimo, en junta convocó al pueblo Aquiles;\\
pues lo inspiró la diosa de blancos brazos, Hera,\\
ansiosa por los dánaos que veía morir.\\
Y luego que, llamados, la junta se reuniera,\\
alzóse el raudo Aquiles, y así empezó a decir:\par
--Atrida, creo que ahora, nuevamente vagando,\\
retroceder debemos, si, la muerte evitando,\\
guerra y peste no rinden juntas a los aqueos.\\
Mas, ea, consultémonos ya de algún adivino\\
sacerdote o intérprete de sueños (pues divino\\
el sueño es también) con que nos diga él sin rodeos\\
qué irrita a Febo Apolo: si nos reclama, acaso,\\
promesas o hecatombes, o si, aceptando el craso\\
humo de albos corderos y de escogidas cabras,\\
\dlgo{nos quitará el plageño.}{Dichas estas palabras,}\\
volvió a sentarse. Entonces se alzó el más excelente\\
augur, Calcas Testórides, quien sabiendo lo que era,\\
es y será, las naves aqueas condujera\\
hasta Ilión, con ayuda del don clarividente\\
que Febo Apolo hiciérale. Y así, benevolente\\
\dlgo{los arengó:}{--Oh Aquiles, caro a Zeus, pues mandarme}\\
quieres que explique la ira de Apolo el real flechero,\\
lo diré; mas prométeme y júrame primero\\
que de palabra y obra socorro has de prestarme;\\
no irrite yo a quien todos los argivos comanda\\
con gran poder, y acatan los aqueos; pues si anda\\
mal con el rey un súbdito como aquél puede más,\\
aunque trague su cólera en el mismo momento\\
hasta satisfacerla guarda el resentimiento\\
en su pecho. Declara, pues, si me salvarás.\par
\dlgo{Contestóle así el rápido Aquiles:}{--Ten confianza,}\\
y el augurio que sabes, dínoslo sin tardanza;\\
pues juro por Apolo, caro a Zeus, y a quien oras,\\
Calcas, cuando a los dánaos su oráculo asesoras,\\
que ni uno de ellos, mientras esté yo vivo y sano,\\
junto a las hondas naves pondrá en ti osada mano;\\
aunque te refieras a Agamenón, que a todos\\
los aqueos se alaba de recuperar ufano.\par
El adivino ilustre tomó confianza y dijo:\\
--No se queja él por votos ni hecatombes, de fijo;\\
mas por su sacerdote que Agamenón vejó,\\
pues no libertó a su hija ni el rescate aceptó.\\
Por esto el Flechero hízonos el mal y todavía\\
nos lo hará, y de las Parcas de la peste sombría\\
no ha de librarnos mientras no quede reparada\\
la ofensa, devolviéndose, sin multa rescatada,\\
al caro padre aquella joven de lindos ojos;\\
y una sacra hecatombe sea en Crisa inmolada.\\
Así conseguiremos aplacar sus enojos.\par
Dichas estas palabras, volvió a sentarse, cuando\\
poseído de angustia se alzó el héroe Atrida,\\
Agamenón potente, la negra entraña henchida\\
de rencor, y con ojos centelleantes mirando\\
\dlgo{a Calcas torvamente, díjole:}{--Mago avieso,}\\
nunca me anuncias nada feliz; antes gozando\\
en predecir desgracias, sólo haces mal con eso.\\
Y ahora entre los dánaos declamas, augurando\\
que el Flechero tortúralos así porque no quise\\
el rescate halagüeño de la doncella Crise,\\
a quien deseaba en casa tener, pues la prefiero\\
realmente a Clitemnestra mi legítima esposa,\\
ya que en cuerpo ni en talle, ni por diestra o graciosa\\
le es inferior; mas, siempre que mejor sea, quiero\\
devolverla, y que el pueblo se salve y no perezca.\\
Pero aprontadme ya otro premio, pues no es cumplido\\
que así el único argivo que ninguno merezca\\
sea yo, viendo todos que el mío habré perdido.\par
Mas, el divino Aquiles respondió de este modo:\\
--Atrida, el más glorioso y avariento de todos,\\
¿qué premios los magnánimos aqueos han de darte,\\
si riquezas comunes no hay en ninguna parte?\\
Repartimos cuanto en las ciudades conquistamos,\\
y no conviene al pueblo que a juntarlo volvamos.\\
Remite al dios la joven, y cuando a los aqueos\\
Zeus nos entregue Troya, la del fuerte bastión,\\
triple o cuádruple pago te asignará el saqueo.\par
Respondió a estas palabras el rey Agamenón:\\
--Por más bravo que seas, no intentes con desvío\\
traerme, ínclito Aquiles, a engaño o persuasión.\\
¿Por conservar tu premio, quieres que deje el mío,\\
y para esto me impones que entre la doncella?\\
Lo haré si los magnánimos aqueos dan por ella\\
algo que la compense, conforme a mi albedrío.\\
Y si no me lo dieran, yo mismo iré a tomar\\
tu premio, o el de Ulises, o el de Ayax, y el que sea\\
se indignará. Mas, de esto, después hemos de hablar.\\
Ahora, pronto botemos a la divina mar\\
negro bajel que una hábil tripulación provea;\\
y una hecatombe, junto con la lozana Crise,\\
embarquemos, poniéndolas al mando de un guerrero:\\
Ayax, Idomeneo, o el deífico Ulises,\\
o tú, Pelida, que eres de todos el más fiero,\\
para que con las víctimas aplaques al Flechero.\par
Mas, de través mirándolo, contestó el raudo Aquiles:\\
--Ay de mí, avaro impúdico ¿habrá un argivo, acaso,\\
que obedezca tus órdenes, ya siguiendo tus pasos,\\
ya afrontando, valiente, los combates viriles?\\
No fue por los lanceros troyanos provocado\\
como vine a la lucha, que en nada me ofendieron,\\
ni jamás mis caballos y vacas se han llevado,\\
ni en Ptía la prolífica y feraz, destruyeron\\
nuestras cosechas, porque siempre nos dividieron\\
muchas sierras umbrías y el mar alborotado.\\
Mas, contigo vinimos, oh infame, a que gozaras\\
con dar en los de Troya venganza a Menelao,\\
y a ti, cara de perro, que ya en nada reparas.\\
Y amenazas privarme del premio que me dieron\\
por mis grandes fatigas los dánaos, aunque poca\\
sale siempre ante el tuyo la parte que me toca\\
de las ricas ciudades troyanas que rindieron.\\
Si bien yo aguanto el peso más arduo de la lucha,\\
y al hacerse el reparto tu recompensa es mucha\\
y la mía pequeña, siempre a mis naves vuelvo\\
contento, aunque cansado de pelear. Mas, resuelvo\\
tornar ya a Ptía, viendo que irme a casa es mejor\\
en las cóncavas naves, que a costa de mi honor,\\
lucro y caudal seguirte granjeando aquí a la vez.\par
\dlgo{Agamenón, rey de hombres, repúsole:}{--Huye, pues,}\\
si se te antoja. No te ruego que a causa mía\\
te quedes. Me harán otros honrosa compañía,\\
y sobre todo el próvido Zeus. Que tú, el más odioso\\
me eres entre los reyes a quienes Dios ayuda,\\
pues de riña y discordias andas siempre gustoso.\\
Y si tu fuerza es tanta, débeslo a un dios, sin duda.\\
Vete con naves y hombres, y reina en paz completa\\
sobre los mirmidones, que tu ira no me inquieta.\\
Mas ya que Febo Apolo va Criseida a quitarme,\\
y con mi buque y gente debo yo enviarla en prenda\\
te amenazo con irme yo en persona a tu tienda,\\
y a la linda Briseida, tu galardón, llevarme;\\
para que veas cómo soy yo el que puede más,\\
y nadie ose conmigo compararse jamás.\par
Tal dijo, y el Peliónida sintió que lo angustiaba\\
su corazón, dudando dentro el pecho velludo,\\
si sacaba a lo largo del muslo el sable agudo\\
y apartando a los otros, al Atrida mataba,\\
o si su furia y cólera reprimía y calmaba.\\
Mientras así agitábanse su corazón y mente,\\
y que desenvainando ya iba la espada ingente,\\
bajó Atenea del cielo, pues la enviaba la diosa\\
de blancos brazos, Hera, que amaba cordialmente\\
a los dos, y por ambos andaba cuidadosa.\\
De atrás tiró al Pelióńida por los cabellos rojos,\\
visible a él solo, que otro ninguno la ha advertido;\\
volvióse aquél atónito, y al punto ha conocido\\
a Palas Atenea por sus terribles ojos;\\
con que aladas palabras así le ha dirigido:\par
--Hija del dios portaégida, ¿a qué viniste? ¿Acaso\\
a presenciar la ofensa de Agamenón Atrida?\\
Mas, puedo predecirte como seguro caso,\\
que pronto su insolencia le costará la vida.\par
Atena la ojizarca contestóle prudente:\\
--He venido del cielo para hacer que, obediente,\\
tu cólera reprimas; pues me envía la diosa\\
de blancos brazos, Hera, que os ama cordialmente\\
a los dos, y que andaba por ambos cuidadosa.\\
Cesa, pues, la disputa, no desnudes la espada,\\
y oféndelo de boca, si hacerlo así te agrada.\\
Te prometo que un día, triple y brillante premio\\
tendrás por este ultraje. Mas, cede a nuestro apremio\\
\dlgo{Y acátanos ya.}{Aquiles de raudos pies contesta:}\\
--Más vale observar, diosa, vuestra orden sin protesta,\\
aunque el corazón tengo grandemente enojado;\\
pues a quien oye a los dioses, es de ellos escuchado.\par
Dijo, y puesta la recia mano al puño argentino,\\
envainó la ancha espada, dócil a Atena, y ella\\
tornó al Olimpo donde con los otros divinos\\
\dlgo{Zeus portaégida mora.}{Mas, con nueva querella,}\\
al Atrida el peliónida increpó siempre acerbo:\par
--Borracho de ojos cínicos y corazón de ciervo,\\
nunca osaste en la guerra con los del pueblo armarte\\
o ir con los principales aqueos al acecho,\\
por miedo a la muerte. Sin duda es más provecho,\\
en el gran campo aqueo despojar de su parte\\
a quien te contradiga. Rey que voraz impera\\
sobre menguada gente, porque de otra manera\\
lanzabas hoy, Atrida, tu último insulto aquí,\\
sobre gran prenda quiero jurarte algo más: ¡Sí,\\
por este cetro que hojas ni ramas dará nunca,\\
pues desde que en el monte dejó su cepa trunca,\\
el bronce al tornearlo privólo, con certeza,\\
de que le retoñaran el follaje ni corteza;\\
y que en sus manos llevan los dánaos al presente,\\
cuando según las leyes de Dios juzgan fielmente:\\
con tal prenda te juro que un día todo aqueo\\
se acordará de Aquiles, y en tu vano deseo\\
de ayudarlos, cuando Héctor tantos mate valiente,\\
desgarrará honda pena en tu corazón furioso,\\
porque al principal de ellos no honraste respetuoso!\par
Tal se expresó el Pelida, contra el suelo arrojando\\
su cetro tachonado de oro, y sentóse enfrente\\
del Atrida que se iba con más furor cargando.\\
Pero entonces de entre ellos se alzó el blandilocuente\\
orador de los pilios, Néstor, cuya cadente\\
palabra, de su lengua, como la miel fluía.\\
Ya dos generaciones de hombres de artuculada\\
voz, que fueron nacidas y consigo criadas\\
en la divina Pylos, ante él pasado habían.\\
Y sobre la tercera mandaba todavía.\\
Este fué el que, arengándolos, benevolente dijo:\par
--¡Oh dioses, qué gran duelo para la tierra aquea!\\
Y cómo a los troyanos, y a Príamo y sus hijos,\\
satisfacción inmensa causaría, de fijo,\\
saber vuestro altercado, siendo que en la pelea\\
y el consejo estáis sobre todos los dánaos. Pero\\
desde que sois más jóvenes, condescended primero;\\
pues yo frecuenté otrora varones principales,\\
mejores que nosotros, y no me desoyeron.\\
Todavía no he visto ni veré hombres iguales\\
a Piritoo, Driante rey de pueblos, Keneo,\\
Exadio y el deífico Polifemo, y Teseo\\
el Egeída, que era como los inmortales.\\
Bravos entre los bravos que alimentó la tierra,\\
ellos a los centauros monteses dieron guerra\\
por más bravos, matándolos en terrible exterminio.\\
Para incorporarme a ellos, dejé la tierra Apía,\\
pues fuí allá desde Pylos, mi lejano dominio,\\
y a su propio llamado luché por cuenta mía.\\
Ningún hombre de ahora resistirlos podría;\\
y, no obstante, acataban mis consejos realmente.\\
Prestadme, así, obediencia, que es lo más conveniente.\\
Ni tú, aunque el mejor seas, la muchacha le quites,\\
puesto que los aqueos con ella lo han premiado.\\
Ni tú, Pelida, quieras altercar, igualado\\
con el rey, ya que en honra ninguno le compite\\
de entre los portacetros por Zeus glorificados.\\
Ahora, si por tu parte tienes más valentía,\\
y si diosa es la excelsa madre que a luz te ha dado,\\
el puede más, pues reina sobre la mayoría.\\
Y tú, Atrida, apacigua tu cólera, y depón\\
tu encono contra Aquiles, que gran resguardo presta\\
a todos los aqueos en la guerra funesta.\par
Contestó, así diciéndole, el rey Agamenón:\\
--Sí, cuanto has dicho, anciano, por cierto es de razón;\\
pero este hombre pretende sobreponerse a todos,\\
y a todos dirigirlos y mandar, de tal modo\\
que nadie ha de aceptárselo. ¿Por que bravo lo hicieron\\
los dioses, que insultara también le permitieron?\par
Mas el divino Aquiles lo interrumpió exclamando:\\
--Por vil y por cobarde con razón pasaría,\\
si a todo cuanto dices cediera. Impón el mando\\
a otros, no a mí, pues nunca más pienso obedecerte.\\
Pero oye esto y consérvalo en tu alma bien grabado:\\
no disputaré a nadie ni a ti, con mano fuerte,\\
la joven, ya que en suma quitáis lo que habéis dado;\\
mas, de lo que me resta junto al negro y ligero\\
bajel, nada podrías tomar, si yo no quiero.\\
De lo contrario, inténtalo; y así éstos sin tardanza\\
verán tu negra sangre chorreando por mi lanza.\par
Habiendo altercado ambos así, en pie se pusieron,\\
y ante la flota aquea la junta disolvieron.\\
El Pelida, a las tiendas y parejos bajeles\\
se fué con Meneciades y sus amigos fieles.\\
Y en tanto, al mar, Atrida, botó un barco ligero;\\
eligió por su propia cuenta veinte remeros,\\
cargó allá una hecatombe para el dios, y llevando\\
a la linda Criseida, la embarcó con esmero.\\
Y el ingenioso Ulises fué quien tomó el comando.\par
Tras que así por la líquida ruta a bogar echaran,\\
mandó Atrida a los pueblos que se purificaran,\\
cual lo hicieron, lanzando toda inmundicia al mar.\\
Luego a Apolo inmolaron en la playa, perfectas\\
hecatombes de toros y de cabras selectas;\\
y el olor iba al cielo con el humo al ondear.\par
Agamenón, en tanto que esto hacía la gente,\\
sin olvidar que a Aquiles en el agrio debate\\
amenazó primero, dijo a sus diligentes\\
heraldos y criados, Taltibio y Euribate:\par
--A la tienda de Aquiles Peleyades id presto,\\
y a la linda Briseida traedme de la mano;\\
o he de ir allá en persona, si él no la da de plano,\\
con más gente a sacársela, y peor le saldrá esto.\par
Despachólos, vertiendo tan fieras expresiones;\\
y ellos, el mar costeando, con ánimo indispuesto,\\
llegaron a las tiendas y barcos mirmidones.\\
Junto a su tienda y negro buque, sentado hallaron\\
a Aquiles, quien, de verlos, no se alegró por cierto;\\
y haciendo al rey la venia, con grave desconcierto,\\
parados ante él nada dijéronle ni hablaron.\\
Pero él entendió en su alma, y alzando así la voz,\\
\dlgo{dijo:}{--Salud, heraldos, mensajeros de Dios}\\
y de los hombres. Aproximaos sin cuidado,\\
pues la culpa no es vuestra, sino de Agamenón,\\
que en busca de la joven Briseida os ha mandado.\\
Vamos, noble Patroclo, saca a disposición\\
de ambos la joven que han de llevarse; y que testigos\\
me sean los dos ante los dioses venturosos,\\
y los hombres mortales, y ese rey rencoroso,\\
si contra el vil desastre cuentan después conmigo.\\
Pues aquél, poseído de furor, ya no sabe\\
ni prever ni acordarse, para que así, al abrigo,\\
combatan los aqueos con él junto a las naves.\par
Dijo, y sacó Patroclo, dócil al caro amigo,\\
de la tienda a la linda Briseida, y ya entregada,\\
tornaron los heraldos, que siguió contrariada\\
la mujer, la aquea flota. Rompiendo en llanto,\\
Aquiles, de los suyos, lejos se apartó en tanto,\\
y sentado a la orilla del espumoso mar,\\
frente al Océano obscuro, con las manos tendidas,\\
a su querida madre comenzó así a implorar:\par
--Madre, ya que me diste tú a luz de corta vida,\\
pudo Zeus, el olímpico altitonante, honrarme.\\
Mas nada hizo, y el déspota Agamenón Atrida,\\
me ha ultrajado, llegando del premio a despojarme.\par
Así exclamó llorando. Lo oyó su augusta madre,\\
que sentada se hallaba junto al anciano padre\\
en el fondo marino, y al punto de la espuma\\
se alzó como una niebla; y ante el lloroso hijo\\
sentóse, acariciándolo con la mano, y le dijo:\par
--Hijo, por qué llorabas? Qué pena tu alma abruma?\\
Habla, y sepámoslo ambos. Nada oculte tu mente.\par
Aquiles contestóle, suspirando hondamente:\\
--Lo sabes. ¿A qué habría de hablar si sabes todo?\\
A Tebas la sagrada ciudad de Eecion fuimos,\\
y habiéndola saqueado, todo acá nos trajimos.\\
Los aqueos partiéronlo de equitativo modo,\\
y eligieron la linda Criseida para Atrida.\\
Mas Crises, sacerdote de Apolo, fué en seguida\\
a las veleras naves aqueas, aportando\\
por rescate de su hija multa inmensa, y llevando\\
en su mano las ínfulas del flechador Apolo\\
a un cetro de oro atadas; con que imploró, no sólo\\
a ambos Atridas, jefes de pueblos, sino a cada\\
\dlgo{uno de los aqueos.}{Fué de éstos aclamado,}\\
respeto al sacerdote, y el rescate aceptado.\\
Mas no plugo al Atrida la idea así acordada,\\
por lo cual, insultándolo, echó a aquél malamente.\\
Fuése el viejo irritado, y Apolo oyó su ruego,\\
que mucho lo quería; con que tiró inclemente\\
saeta a los argivos, cuyos soldados luego\\
morían a montones; pues las divinas flechas,\\
por todo el grande ejército Aqueo hacían brechas.\\
Díjonos los oráculos del Flechero un sapiente\\
adivino, y consejo di yo inmediatamente\\
de apaciguar al numen. Entonces el Atrida,\\
encendiéndose en cólera, alzóse bruscamente,\\
y me echó una amenaza que está, en verdad, cumplida.\\
Pues ya los perspicaces dánaos a Crise enviaron\\
junto con reales dones, en nave bogadora;\\
mientras que mi tienda unos heraldos se llevaron\\
a la hija de Briseo que aquéllos me otorgaron.\\
Si a tu hijo auxiliar quieres, ve al Olimpo e implora\\
a Zeus, que acaso un día su corazón tocaste\\
con hechos o palabras; porque recuerdo ahora,\\
que en el paterno alcázar mil veces te gloriaste\\
de que sola entre todos los inmortales fueras\\
quien al Kronión sombrío vergonzoso contraste\\
supo evitar, cuando otros Olímpicos, con Hera\\
y Poseidón, y Palas Atenea, qusieran\\
encadenarlo. Pero tú, diosa, allí acudiste,\\
desataste sus lazos, y al punto condujiste\\
al vasto Olimpo el fiero centímano que llaman\\
Briareo los dioses y Egeón los humanos.\\
Más fuerte que su padre, fué, de su gloria ufano,\\
a sentarse a la vera del Kronión; y medrosos\\
de él, ya al otro no ataron los dioses venturosos.\\
Recuérdaselo ahora, y llégate a él, y abrazando\\
sus rodillas, consigue que a los troyanos dando\\
su ayuda, ante la flota rechacen y destrocen\\
estos a los aqueos, que así de su rey gocen.\\
Y su error vea el déspota Agamenón Atrida,\\
porque al principal de ellos no honró en forma cumplida.\par
\dlgo{Respondióle llorando Tetis:}{--¿Por qué, hijo mío,}\\
si te di a luz en hora fatal, te habré criado?\\
Sin lágrimas ni pena debiste en los navíos\\
quedar, ya que a tan próximo fin estás destinado.\\
Y ahora eres corto en días y el más desventurado;\\
que al ser en el palacio te di con hado impío.\\
A hablar por ti al fulmíneo Zeus iré yo al nevado\\
Olimpo, y quizá pueda persuadirlo. Y ahora,\\
con los dánaos airado, gana tu bogadora\\
flota, y no hagas más guerras. Hacia el Océano ayer\\
partió Zeus a un banquete que dan los caballeros\\
etíopes y todos los dioses con él fueron.\\
Mas al Olimpo en doce días debe volver;\\
y a su alcázar de bronce, yendo en tu nombre a hablarle,\\
asiré sus rodillas y espero propiciarle.\par
Dijo y partió, dejándolo con el alma enconada,\\
por la elegante esposa que le fué arrebatada.\\
Y mientras tanto Ulises a Crisa conducía\\
la hecatombe sagrada. Tan luego que arribaron\\
al puerto de hondas aguas, amainando guardaron\\
las velas en el negro buque; y en la crujía,\\
al mástil, prontamente, con cables abatieron.\\
Bogaron a remo hasta la rada donde echaron\\
anclas; y atando amarras, a la costa bajaron.\\
sacaron la hecatomble para Apolo el flechero,\\
y al fin salió Criseida del bajel marinero.\\
El ingenioso Ulises condújola al altar,\\
y entregándola al caro padre empezó así a hablar:\par
--Oh Crises, el rey de hombres Agamenón me envía\\
a que tu hija te traiga, y a que en ofrenda pía\\
inmole yo a aquí a Febo la hecatombe sagrada\\
en nombre de los dánaos, a ver si el dios se apiada,\\
pues mandó a los argivos lamentable avería.\par
Dijo, y su hija entrególe, que él recibió contento.\\
Y al punto, en torno al ara de sólido cimiento,\\
para el dios ordenaron la hecatombe sagrada.\\
Laváronse las manos, tomaron la salada\\
harina, y allá entonces, en alta voz orando,\\
Crises rogó por ellos, las manos levantando:\par
--Oyeme, dios del areo de plata, protector\\
de Crisa y de la santa Kila, dominador\\
de Ténedos: si un día mis ruegos escuchando\\
me honraste, al pueblo aqueo cruelmente castigando,\\
ahora como entonces, cúmpleme también este voto,\\
y libra a los dánaos de la horrorosa peste.\par
Dijo así, y fué de Apolo su oración escuchada.\\
Luego que oran y esparcen ya la harina salada,\\
mancornando a las víctimas, degüellan y desuellan.\\
Cortan los muslos, cúbrenlos con un doble haz de pella\\
y magras, y el anciano sobre la lecha hacheada\\
quemólos, y con negro vino los fué rociando.\\
En torno suyo había mancebos que empuñando\\
asadores de cinco puntas, se dispusieron.\\
Consumidos los muslos, las entrañas comieron;\\
y ensartando en los pinchos lo restante, lo asaron\\
con cuidado, por trozos; sacáronlo, y cumplida\\
tal faena, el dispuesto banquete comenzaron\\
sin que a nadie faltara su porción. En seguida\\
que saciaron sed y hambre, los mancebos colmaron\\
las cráteras de vino que después repartían,\\
y a todos las primicias de la copa obsequiaron.\\
Los jóvenes aqueos, tras de holgar todo el día,\\
al dios se propiciaron cantándole un hermoso\\
himno, que oyó el Flechero con ánimo gozoso.\par
Cuando el sol se entró y vino la sombra, se acostaron\\
junto al cable del buque; y al asomar temprana\\
la Aurora rosodáctila, hija de la mañana,\\
para el campo aqueo, mar afuera tomaron.\\
Y el flechador Apolo buen viento les envía.\\
Izando luego el mástil, las velas desplegaron,\\
que en blanca combadura con el soplo se hincharon,\\
mientras la honda purpúrea resonando batía\\
la quilla del navío que su rumbo corría.\\
Ya ante el gran campo aqueo, la nave a tierra halaron;\\
y empinándola en grandes vigas sobre la arena,\\
por los buques y tiendas luego se dispersaron.\par
Junto a su flota en tanto, con su alma de ira llena,\\
seguía el raudo Aquiles, noble hijo de Peleo;\\
y de la ilustre junta y el combate alejado,\\
el corazón ya le iba devorando el deseo\\
\dlgo{de griterió y Guerra.}{Mas, cuando hubo llegado}\\
la duodécima aurora, y al Olimpo volvieron\\
los sempiternos dioses por Zeus encabezados,\\
Tetis, que los encargos de su hijo no ha olvidado,\\
emergiendo temprana de las olas del mar,\\
y cruzando el gran cielo, fué en el Olimpo a hallar\\
al vidente Krónida, que se hallaba sentado\\
en el pico más alto del Olimpo escarpado.\\
Fué, pues, y ante él sentándose, las rodillas le asió\\
con la izquierda, y tomándole el mentón con la diestra,\\
su ruego al soberano dios Kronión dirigió:\par
--Padre Zeus, si he dado entre los inmortales muestra\\
de serte útil con actos o palabras, te pido\\
que oigas mi ruego y honres a mi hijo, el de más corta\\
vida entre todos, porque de Agamenón soporta\\
la ofensa que arrancándole su premio le ha inferido.\\
Véngalo, pues, Olímpico Zeus próvido, acordando\\
el triunfo a los troyanos, hasta que los aqueos\\
cumplan con mi hijo y lo honren conforme a sus deseos.\par
Dijo, mas Zeus nubígero, largo tiempo callando,\\
nada respondió. Tetis, que abrazando seguía\\
sus rodillas con fuerza, suplicó todavía:\par
--Prométeme y con una señal veraz asiente,\\
o niega, ya que a nadie temes; y así enterada,\\
sabré si entre las diosas soy la más despreciada.\par
El nubígero Zeus, suspirando hondamente,\\
\dlgo{díjole:}{--Algo funesto se apronta, pues con Hera}\\
me enconarás, cuando ella me provoque y zahiera.\\
Ya entre los dioses ríñeme sin motivo, a pretexto\\
de que yo en la pelea protejo a los troyanos.\\
Pero vete y apártate, no te advierta Hera; y a esto,\\
para que bien se cumpla, déjalo ya en mis manos.\\
Te haré con la cabeza la seña en que confías.\\
Entre los inmortales no hay mayor garantía.\\
Que irrevocable y cierto, de cumplirse no deja\\
lo que con la cabeza ratifiqué a fe mía.\par
Dijo el Kronión. Al signo de las azules cejas,\\
ondeó el fragante pelo del rey en la inmortal\\
cabeza, y el Olimpo retembló colosal.\par
Dicho aquello, apartáronse. Regresó ella a su centro,\\
al hondo mar saltando del Olimpo esplendente,\\
y Zeus a su palacio. Los dioses juntamente\\
ante su padre alzáronse, saliéndole al encuentro.\\
Sentóse él en su trono; mas Hera que con tino\\
había visto cómo con él deliberó\\
Tetis pies de plata, hija del anciano marino,\\
al dios Kronión hirientes palabras dirigió:\par
--¿Cuál de los dioses, pérfido, contigo a conversado?\\
Siempre tramar secretos sin mí, te causa agrado;\\
mas, de lo que meditas, nunca amable me enteras.\par
\dlgo{Contestó el padre de hombres y dioses:}{--En vano, Hera,}\\
pretendes de tal modo saber todas mis cosas,\\
porque arduas te saldrían aunque seas mi esposa.\\
Lo que es de oírse, ni hombre ni dios sabrá primero;\\
mas todo cuanto sin los dioses resolver quiero,\\
\dlgo{No indagues ni preguntes.}{Respondió la augusta Hera}\\
\dlgo{de ojos de buey:}{--¡Terrible Kronida, qué dijiste!}\\
Poco indago o pregunto, puesto que deliberas\\
muy tranquilo tus cosas. Mas tengo el alma triste\\
de recelar que acaso te sedujo con tino\\
Tetis pies dep plata, hija del anciano marino.\\
Temprano vino a asirte las rodillas, y creo\\
que para honrar a Aquiles le acordaste, de fijo,\\
perder a muchos ante los bajeles aqueos.\par
El nubígero Zeus, respondiéndole dijo:\\
--Desgraciada! Sospechas siempre lo que no oculto.\\
Así obtendrás tan sólo quedar más alejada\\
de mi afecto, y con ello sufrir mayor insulto.\\
Si es verdad lo que dices, será porque me agrada.\\
Pero, acata mis órdenes y siéntate callada.\\
Teme que ni los dioses del Olimpi, a ti amables,\\
te valgan si en ti pongo mis manos formidables.\par
Dijo así. La augusta Hera de ojos de buey, temiendo,\\
fué a sentarse callada, su furor reprimiendo.\\
Y gimieron los célicos dioses en la morada\\
de Zeus. Mas ya el artífice Hefesto los arenga\\
mimando a Hera de blancos brazos, su madre amada:\par
--Temo que algo insufrible y aciago sobrevenga,\\
si así por los mortales reñís, alborotando\\
a los dioses, porque hasta se nos irá amargando\\
el placer de las fiestas. Yo aconsejo a mi madre,\\
aun cuando es tan prudente, mime a Zeus, caro padre;\\
con que así él nuestras fiestas no turbe rencilloso.\\
Pues si viene al Olímpico fulminador la idea,\\
puede echarnos abajo, siendo el más poderoso.\\
Hazle en dulces palabras halago cariñoso,\\
a fin de que el Olímpico favorable nos sea.\par
Habló así; y dirigiéndose a su madre querida,\\
puso en sus manos una doble copa y le dijo:\\
--Sufre todo, madre, aunque sea cruel tu quebranto.\\
No te vea golpeada yo, queriéndote tanto;\\
y por más que en su furia nada te valga tu hijo,\\
pues arduo es al Olímpico resistir. Ya una vez\\
que intenté socorrerte, me asió de un pie y después\\
desde el umbral divino me arrojó. Todo el día\\
rodé, cayendo en Lemnos cuando el sol se ponía.\\
Débil soplo quedábame de vida postrimera,\\
y los sintios salváronme del golpe prontamente.\par
Así dijo, y la diosa de blancos brazos, Hera,\\
la copa de su mano recibió sonriente.\\
Luego él, por la derecha sirvió a los otros dioses,\\
sacando de una crátera el dulce néctar. A esto,\\
de los felices númenes inextinguibles alzóse\\
la risa en el palacio, viendo afanarse a Hefesto\par
Hasta que el sol se puso, todo el día gozaron\\
del festín, sin que nadie justa porción faltara,\\
ni la ira magnífica de Apolo, ni la clara\\
voz con que allá las musas alternando cantaron.\par
Mas, luego de apagada la solar refulgencia,\\
ganaron las mansiones que a cada uno había hecho\\
Hefesto el cojo ilustre con ingeniosa ciencia.\\
Zeus, el fulmíneo Olímpico, fué a buscar en su lecho\\
la calma y dulce sueño que hallaba satisfecho.\\
Habiendo allá subido, durmióse; y a su vera,\\
reposaba acostada la crisótrona Hera.
\pend
\endnumbering